%% LyX 2.5.1 created this file.  For more info, see https://www.lyx.org/.
%% Do not edit unless you really know what you are doing.
\documentclass[12pt,american]{extarticle}
\usepackage[T1]{fontenc}
\usepackage[utf8]{inputenc}
\usepackage{mathtools}
\usepackage{amsmath}
\usepackage{amsthm}
\usepackage{amssymb}
\usepackage{geometry}
\geometry{verbose,tmargin=2.5cm,bmargin=2.5cm,lmargin=2cm,rmargin=3cm}
\usepackage{setspace}
\onehalfspacing

\makeatletter
%%%%%%%%%%%%%%%%%%%%%%%%%%%%%% Textclass specific LaTeX commands.
\numberwithin{equation}{section}
\numberwithin{figure}{section}

%%%%%%%%%%%%%%%%%%%%%%%%%%%%%% User specified LaTeX commands.
\usepackage[utf8]{inputenc}
\usepackage[T1]{fontenc}
% \usepackage[default]{comicneue}

\usepackage{hyperref}
\hypersetup{
	colorlinks=true, %set true if you want colored links
	linktoc=all,     %set to all if you want both sections and subsections linked
	linkcolor=black,  %choose some color if you want links to stand out
}

% \usepackage{url}

%\usepackage[style=alphabetic]{biblatex}
%\addbibresource{bibliography.bib}

\usepackage{graphicx}
\usepackage{enumitem}
\usepackage{booktabs}
\usepackage{csquotes}
% \usepackage{emptypage}
\usepackage{subcaption}
\usepackage{multicol}
\usepackage[dvipsnames]{xcolor}
% \usepackage[usenames, dvipsnames]{xcolor}
\usepackage{amsmath, mathtools, amssymb}
% \usepackage{amsfonts} 
\usepackage{newpxtext} 
\usepackage[euler-digits]{eulervm}
%\usepackage{libertine,libertinust1math}

%\usepackage[sc,osf]{mathpazo}   % With old-style figures and real smallcaps.
%\linespread{1.025}              % Palatino leads a little more leading

% Euler for math and numbers
%\usepackage[euler-digits,small]{eulervm}
\usepackage{extarrows}
\usepackage{amsthm}
\usepackage{mathrsfs}
\usepackage{cancel}
\usepackage{bm}
\usepackage{systeme}
\usepackage{stmaryrd}
\usepackage{dsfont}


\let\implies\Rightarrow
\let\impliedby\Leftarrow
\let\iff\Leftrightarrow
\let\epsilon\varepsilon
\newcommand\contra{\scalebox{1.1}{$\lightning$}}
\let\emptyset\varnothing


\definecolor{correct}{HTML}{009900}
\newcommand\correct[2]{\ensuremath{\:}{\color{red}{#1}}\ensuremath{\to }{\color{correct}{#2}}\ensuremath{\:}}
\newcommand\green[1]{{\color{correct}{#1}}}


\newcommand\hr{
	\noindent\rule[0.5ex]{\linewidth}{0.5pt}
}


% hide parts
\newcommand\hide[1]{}



% si unitx
\usepackage{siunitx}
\sisetup{locale = FR}
% \renewcommand\vec[1]{\mathbf{#1}}
\newcommand\mat[1]{\mathbf{#1}}


% tikz
\usepackage{tikz}
\usepackage{tikz-cd}
\usetikzlibrary{intersections, angles, quotes, calc, positioning}
\usetikzlibrary{arrows.meta}
\usepackage{pgfplots}
\pgfplotsset{compat=1.13}



\tikzset{
	force/.style={thick, {Circle[length=2pt]}-stealth, shorten <=-1pt}
}

% theorems
\makeatother
\usepackage{thmtools}
\usepackage[framemethod=TikZ]{mdframed}
\mdfsetup{skipabove=1em,skipbelow=0em}


% \theoremstyle{plain}

\declaretheoremstyle[
	headfont=\bfseries\sffamily\color{ForestGreen!70!black}, bodyfont=\normalfont,
	mdframed={
			linewidth=2pt,
			rightline=false, topline=false, bottomline=false,
			linecolor=ForestGreen, backgroundcolor=ForestGreen!5,
		}
]{thmgreenbox}

\declaretheoremstyle[
	headfont=\bfseries\sffamily\color{NavyBlue!70!black}, bodyfont=\normalfont,
	mdframed={
			linewidth=2pt,
			rightline=false, topline=false, bottomline=false,
			linecolor=NavyBlue, backgroundcolor=NavyBlue!5,
		}
]{thmbluebox}

\declaretheoremstyle[
	headfont=\bfseries\sffamily\color{RoyalPurple}, bodyfont=\normalfont\itshape,
	mdframed={
			linewidth=2pt,
			rightline=false, topline=false, bottomline=false,
			linecolor=Rhodamine, backgroundcolor=Rhodamine!5,
		}
]{thmpinkbox}

\declaretheoremstyle[
	headfont=\bfseries\sffamily\color{NavyBlue!70!black}, bodyfont=\normalfont,
	mdframed={
			linewidth=2pt,
			rightline=false, topline=false, bottomline=false,
			linecolor=NavyBlue
		}
]{thmblueline}

\declaretheoremstyle[
	headfont=\bfseries\sffamily\color{RawSienna!70!black}, bodyfont=\normalfont\itshape,
	mdframed={
			linewidth=2pt,
			rightline=false, topline=false, bottomline=false,
			linecolor=RawSienna, backgroundcolor=RawSienna!5,
		}
]{thmredbox}

\declaretheoremstyle[
	headfont=\bfseries\sffamily\color{RawSienna!70!black}, bodyfont=\normalfont,
	numbered=no,
	mdframed={
			linewidth=2pt,
			rightline=false, topline=false, bottomline=false,
			linecolor=RawSienna, backgroundcolor=RawSienna!1,
		},
	qed=\qedsymbol
]{thmproofbox}

\declaretheoremstyle[
	headfont=\bfseries\sffamily\color{NavyBlue!70!black}, bodyfont=\normalfont,
	numbered=no,
	mdframed={
			linewidth=2pt,
			rightline=false, topline=false, bottomline=false,
			linecolor=NavyBlue, backgroundcolor=NavyBlue!1,
		},
]{thmexplanationbox}

\declaretheorem[style=thmgreenbox, name=Definition, numberwithin=section]{definition}
\declaretheorem[style=thmbluebox, name=Example, numberwithin=section]{eg}
\declaretheorem[style=thmpinkbox, name=Exercise, numberwithin=section]{ex}
\declaretheorem[style=thmredbox, name=Proposition, numberwithin=section]{prop}
\declaretheorem[style=thmredbox, name=Theorem, numberwithin=section]{theorem}
\declaretheorem[style=thmredbox, name=Lemma, numberwithin=section]{lemma}
\declaretheorem[style=thmredbox, numbered=no, name=Corollary, numberwithin=section]{corollary}

\declaretheorem[style=thmproofbox, name=Proof]{replacementproof}
\renewenvironment{proof}[1][\proofname]{\vspace{-10pt}\begin{replacementproof}}{\end{replacementproof}}


\declaretheorem[style=thmexplanationbox, name=Proof]{tmpexplanation}
\newenvironment{explanation}[1][]{\vspace{-10pt}\begin{tmpexplanation}}{\end{tmpexplanation}}

\declaretheorem[style=thmblueline, numbered=no, name=Remark]{remark}
\declaretheorem[style=thmblueline, numbered=no, name=Note]{note}

\newtheorem*{uovt}{UOVT}
\newtheorem*{notation}{Notation}
\newtheorem*{previouslyseen}{As previously seen}
\newtheorem*{problem}{Problem}
\newtheorem*{observe}{Observe}
\newtheorem*{property}{Property}
\newtheorem*{intuition}{Intuition}


\usepackage{etoolbox}
\AtEndEnvironment{vb}{\null\hfill$\diamond$}%
\AtEndEnvironment{intermezzo}{\null\hfill$\diamond$}%
% \AtEndEnvironment{opmerking}{\null\hfill$\diamond$}%

% http://tex.stackexchange.com/questions/22119/how-can-i-change-the-spacing-before-theorems-with-amsthm
\makeatletter
% \def\thm@space@setup{%
%   \thm@preskip=\parskip \thm@postskip=0pt
% }

\newcommand{\oefening}[1]{%
	\def\@oefening{#1}%
	\subsection*{Oefening #1}
}

\newcommand{\suboefening}[1]{%
	\subsubsection*{Oefening \@oefening.#1}
}

\newcommand{\exercise}[1]{%
	\def\@exercise{#1}%
	\subsection*{Exercise #1}
}

\newcommand{\subexercise}[1]{%
	\subsubsection*{Exercise \@exercise.#1}
}


\usepackage{xifthen}

\def\testdateparts#1{\dateparts#1\relax}
\def\dateparts#1 #2 #3 #4 #5\relax{
	\marginpar{\small\textsf{\mbox{#1 #2 #3 #5}}}
}

\def\@lesson{}%
\newcommand{\lesson}[3]{
	\ifthenelse{\isempty{#3}}{%
		\def\@lesson{Lecture #1}%
	}{%
		\def\@lesson{Lecture #1: #3}%
	}%
	\subsection*{\@lesson}
	\testdateparts{#2}
}

% \renewcommand\date[1]{\marginpar{#1}}


% fancy headers

\usepackage{fancyhdr}
\pagestyle{fancy}
% \fancyhead[LE,RO]{Gilles Castel}
% \fancyhead[RO,LE]{\@lesson}
% \fancyhead[RE,LO]{}
% \fancyfoot[LE,RO]{\thepage}
% \fancyfoot[C]{\leftmark}
% \fancyhf{}
\fancyhf[roh,leh]{\thepage}
\fancyhf[reh]{\nouppercase{\leftmark}}
\fancyhf[loh]{\nouppercase{\rightmark}}
\cfoot{}
\setlength{\headheight}{15pt}
\renewcommand{\headrulewidth}{0pt}
\renewcommand{\footrulewidth}{0pt}


\makeatother




% notes
% \usepackage{todonotes}
\usepackage{tcolorbox}

\tcbuselibrary{breakable}
\newenvironment{verbetering}{\begin{tcolorbox}[
			arc=0mm,
			colback=white,
			colframe=green!60!black,
			title=Opmerking,
			fonttitle=\sffamily,
			breakable
		]}{\end{tcolorbox}}

\newenvironment{noot}[1]{\begin{tcolorbox}[
			arc=0mm,
			colback=white,
			colframe=white!60!black,
			title=#1,
			fonttitle=\sffamily,
			breakable
		]}{\end{tcolorbox}}


\newcommand{\outermeasure}{{{\mu^{\star}}}}

%\def\mbb#1{\mathbb{#1}}
%\def\mfk#1{\mathfrak{#1}}

\newcommand{\inv}{^{\raisebox{.2ex}{$\scriptscriptstyle-1$}}}

\makeatother

\theoremstyle{definition}
\newtheorem*{xca*}{\protect\exercisename}
\newtheorem*{sol*}{\protect\solutionname}
\usepackage{babel}
\providecommand{\exercisename}{Exercise}
\providecommand{\solutionname}{Solution}

\begin{document}
\title{Aluffi 1-5 solutions}
\author{pika}

\maketitle
\global\long\def\bN{\mathbb{N}}%

\global\long\def\bQ{\mathbb{Q}}%

\global\long\def\bZ{\mathbb{Z}}%

\global\long\def\bR{\mathbb{R}}%

\global\long\def\bC{\mathbb{C}}%

\global\long\def\bB{\mathbb{B}}%

\global\long\def\powerset{\mathscr{P}}%

\global\long\def\image{\operatorname\{Im\}}%

\global\long\def\dom{\operatorname{dom}}%

\global\long\def\id{\operatorname\{Id\}}%

\global\long\def\Span{\operatorname\{span\}}%

\global\long\def\one{\mathds{1}}%

\global\long\def\calA{\mathcal{A}}%

\global\long\def\calB{\mathcal{B}}%

\global\long\def\sigmaalgebra{\mathcal{M}}%

\global\long\def\calR{\mathcal{R}}%

\global\long\def\calE{\mathcal{E}}%

\global\long\def\Hom{\mathsf{Hom}}%

\global\long\def\End{\mathsf{End}}%

\global\long\def\Aut{\mathsf{Aut}}%

\global\long\def\catC{\mathsf{C}}%

\global\long\def\Obj{\mathsf{\mathsf{Obj}}}%

\global\long\def\morphism{\longrightarrow}%


\begin{xca*}
\cite[Ch-1, Ex-5.1]{aluffi}\label{aluffi51} Prove that a final object
in a category $\catC$ is initial in the opposite category $\catC^{\mathsf{op}}$. 
\end{xca*}
\begin{sol*}
Suppose $I$ is the final object in $\catC.$ Then, take any object
$Z\in\catC^{\mathsf{op}}.$ Since there is always a unique morphism
$Z\to I$ in $\catC$, there is a unique morphism $I\to Z$ in $\catC^{\mathsf{op}}.$
Siince $Z\in\mathsf{C}^{\mathsf{op}}$ was arbitrary, this proves
that $I$ is the initial object in $\catC^{\mathsf{op}}$. $\hfill\blacksquare$
\end{sol*}
\begin{xca*}
\cite[Ch-1, Ex-5.2]{aluffi} Prove that $\emptyset$ is the \emph{unique
}initial object in $\mathsf{Set}$. 
\end{xca*}
\begin{sol*}
Suppose $X\in\Obj(\mathsf{Set})$ is an initial object. Then, there
is a unique morphism $X\to\emptyset$ which proves that $X=\emptyset.$$\hfill\blacksquare$
\end{sol*}
\begin{xca*}
\cite[Ch-1, Ex-5.3]{aluffi} Prove that final objects are unique up
to isomorphisms. 
\end{xca*}
\begin{sol*}
Suppose there are two final objects $F_{1},F_{2}$ in a category $\catC.$
Then, since $F_{2}$ is final, we pick the unique morphism $f:F_{1}\to F_{2}$
and since $F_{1}$ is final we consider the unique morphism $g:F_{2}\to F_{1}$.
Then, notice that $fg=1_{F_{2}}$. This is because $F_{2}$ is final,
so there is only a unique morphism $1_{F_{2}}:F_{2}\to F_{2}$ and
$gf$ must be identical to this. By a similar argument we deduce that
$gf=1_{F_{1}}.$ This completes the proof that any two final objects
$F_{1}$ and $F_{2}$ are always isomorphic. $\hfill\blacksquare$
\end{sol*}
\begin{xca*}
\cite[Ch-1, Ex-5.4]{aluffi} What are the final objects in the category
of 'pointed sets' in Eg-3.8? Are they unique? 
\end{xca*}
\begin{sol*}
Any singleton $\{*\}$ with the point $*$ is a final object. Indeed,
given any object $(Z,z)$ in $\mathsf{Set}^{*}$ there is obviusly
a morphism $\sigma:Z\to\{*\}$ so that $\sigma(z)=*$. This is clearly
the constant map $Z\to\{*\}$. Any other singleton set $S$ with its
point $s$ forming the pair $(S,s)$ would be a final object in this
category so such a construction is not unique although obviosuly by
a previous exercise it would be unique up to isomorphism. $\hfill\blacksquare$
\end{sol*}
\begin{xca*}
\cite[Ch-1, Ex-5.5]{aluffi} What are the final objects in the category
considered in Eg-5.3? 
\end{xca*}
\begin{sol*}
The final object should be $(c:a\mapsto*,\{*\})$. Indeed, in this
way, given any object $(f,Z)$ where $f$ agrees on equivalent points
of $A$, we get a morphism $\sigma:(f,Z)\to(c,\{*\})$ that is constant
in $Z$ (maps every point in $Z$ to $*$). Such a morphism is guaranteed
to be unique. $\hfill\blacksquare$
\end{sol*}
\begin{xca*}
\cite[Ch-1, Ex-5.6]{aluffi} Consider the category corresponding to
Eg-3.3 by endowing the set $\bZ_{>0}$ of positive integers with the
\emph{divisibility relation. }Thus, there is exactly one morphism
$d\to m$ in this category if and only if $d$ divides $m$ without
remainder; there is no morphism between $d$ and $m$ otherwise. Show
that this category has products and coproducts. What are their conventional
names? 
\end{xca*}
\begin{sol*}
That this construction defines a category is obvious so we will not
bother with proving that part. Notice that the $\gcd(a,b)$ divides
both $a$ and $b$ so there are morphisms $\gcd(a,b)\to a,\gcd(a,b)\to b$.
Next, given any other integer $z$ with pairs of morphisms $z\to a$
and $z\to b$, there is a (unique by construction of this category)
morhpism $z\to\gcd(a,b)$ simply by definition of $\gcd$(a,b). $\hfill\blacksquare$
\end{sol*}
\begin{xca*}
\cite[Ch-1, Ex-5.8]{aluffi} Show that in every category $\catC$
the products $A\times B$ and $B\times A$ are isomorphic, if they
exist. 
\end{xca*}
\begin{sol*}
Suppose $\catC$ is a category with objects $A,B$ so that the products
$A\times B$ and $B\times A$ both exist. We denote the canonical
projections by $p_{A}:A\times B\to A,p_{B}:A\times B\to B,p'_{B}:B\times A\to B,p'_{A}:B\times A\to A$.
Next, by the universal property of $A\times B$, we get a unique morphism
$f:B\times A\to A\times B$ so that $p_{A}f=p'_{A},p_{B}f=p'_{B}.$
Next, by the universal property of $B\times A,$ there is a unique
morphism $g:A\times B\to B\times A$ so that 
\[
p'_{B}g=p_{B},p'_{A}g=p_{A}
\]

Now, we see that $p_{A}f=p'_{A}gf=p'_{A}$ and $p_{B}f=p'_{B}gf=p'_{B}.$
Also, we see that $p'_{B}g=p_{B}fg=p_{B}$ and $p'_{A}g=p_{A}fg=p_{A}.$
Notice that $1_{B\times A}$ is a morphism $B\times A\to B\times A$
that satisfies the first two relations the same way as $gf$. By uniqueness,
we get that $gf=1_{B\times A}$. Similarly, $1_{A\times B}$ is a
morphism $A\times B\to A\times B$ that satisfies the third and the
fourth commutativity relations above the same way as $fg$. Once again,
by uniqueness we find that $fg=1_{A\times B}.$ This completes the
proof. $\hfill\blacksquare$ 
\end{sol*}
\begin{xca*}
\cite[Ch-1, Ex-5.9]{aluffi} Let $\catC$ be a category with products.
Find a reasonable candidate for the universal property that the product
$A\times B\times C$ of three objects of $\catC$ ought to satisfy,
and prove that both $(A\times B)\times C$ and $A\times(B\times C)$
satisfy this universal property. Deduce that $(A\times B)\times C$
and $A\times(B\times C)$ are necessarily isomorphic. 
\end{xca*}
\begin{sol*}
In any category $\catC$ in which binary products are defined, there
exists an object $P\coloneqq\prod^{3}_{i=1}A_{i}$ and morphisms $(\pi_{i}:P\to A_{i})^{3}_{1}$
so that for any object $Z\in\Obj(\catC)$ and family of morphisms
$(\varphi_{i}:Z\to A_{i})^{3}_{1}$ there exists a uniquely defined
morphism $\varphi:Z\to P$ so that the diagram below commutes for
every $1\leq i\leq3$. 
\end{sol*}
% https://q.uiver.app/#q=WzAsMyxbMiwwLCJaIl0sWzIsMiwiUCJdLFswLDIsIkFfaSAiXSxbMCwyLCJcXHZhcnBoaV9pIiwyXSxbMSwyLCJcXHBpX2kiXSxbMCwxLCJcXHZhcnBoaSJdXQ==
\[\begin{tikzcd}
	&& Z \\
	\\
	{A_i } && P
	\arrow["{\varphi_i}"', from=1-3, to=3-1]
	\arrow["\varphi", from=1-3, to=3-3]
	\arrow["{\pi_i}", from=3-3, to=3-1]
\end{tikzcd}\]

There are three such commutative diagrams for $i=1,2,3$ where $A_{1}=A,A_{2}=B,A_{3}=C$.
The commutativity relations are as follows

\[
\pi_{i}\circ\varphi=\varphi_{i},\qquad1\leq i\leq3.
\]

We will only show that $(A_{1}\times A_{2})\times A_{3}$ satisfies
this universal property. Indeed, notice that there are projections
$\pi':(A_{1}\times A_{2})\times A_{3}\to(A_{1}\times A_{2})$ and
$\pi_{3}:(A_{1}\times A_{2})\times A_{3}\to A_{3}$ Then, suppose
$Z$ is any other object with morphisms $(\varphi_{i}:Z\to A_{i})^{3}_{1}$.
Next, notice that there is a unique morhpism $\varphi_{1}\times\varphi_{2}:Z\to A_{1}\times A_{2}$
by the universal property of the product $A_{1}\times A_{2}$ which
satisfies 
\[
\pi'_{1}\circ(\varphi_{1}\times\varphi_{2})=\varphi_{1},
\]
\[
\pi'_{2}\circ(\varphi_{1}\times\varphi_{2})=\varphi_{2}
\]

where $\pi'_{1}:(A_{1}\times A_{2})\to A_{1}$ and $\pi'_{2}:(A_{1}\times A_{2})\to A_{2}$
are the canonical projections of the product $A_{1}\times A_{2}$
to its factors $A_{1}$ and $A_{2}$ respectively.

Next, due the the universal property of the binary product $(A_{1}\times A_{2})\times A_{3}$,
there is a unique morphism $(\varphi_{1}\times\varphi_{2})\times\varphi_{3}:Z\to(A\times B)\times C$
so that 

\begin{align*}
\pi'\circ((\varphi_{1}\times\varphi_{2})\times\varphi_{3})= & \varphi_{1}\times\varphi_{2},\\
\pi_{3}\circ((\varphi_{1}\times\varphi_{2})\times\varphi_{3})= & \varphi_{3}
\end{align*}

We notice that $(A_{1}\times A_{2})\times A_{3}$ has projections
to $A_{1}$ as $\pi_{1}\coloneqq\pi'_{1}\circ\pi'$, to $A_{2}$ as
$\pi_{2}\coloneqq\pi'_{2}\circ\pi'$ and to $A_{3}$ as $\pi_{3}$.
This tells us that 
\begin{align*}
\pi_{1}\circ((\varphi_{1}\times\varphi_{2})\times\varphi_{3})= & \pi'_{1}\circ(\varphi_{1}\times\varphi_{2})\\
= & \varphi_{1}
\end{align*}

and 
\begin{align*}
\pi_{2}\circ((\varphi_{1}\times\varphi_{2})\times\varphi_{3}) & =\pi'_{2}\circ(\varphi_{1}\times\varphi_{2})\\
 & =\varphi_{2}
\end{align*}

and the final equality is
\[
\pi_{3}\circ((\varphi_{1}\times\varphi_{2})\times\varphi_{3})=\varphi_{3}
\]

We simply define our $\varphi$ to be this morphism $(\varphi_{1}\times\varphi_{2})\times\varphi_{3}$
and clearly from the above proof it is uniquely defined. 

This completes the proof that $(A_{1}\times A_{2})\times A_{3}$ also
satisfies the universal property of products of triplets. 

For uniqueness, suppose $Z$ with the morphisms $(\rho_{i}:Z\to A_{i})^{3}_{1}$
also satisfies the same universal property. Then, by universal property
of $P$ there is a unique morphism $\varphi:Z\to P$ that satisfies
$\pi_{i}\circ\varphi=\rho_{i}$ for every $1\leq i\leq3$. Again,
by universal property of $Z$ there is a unique morphism $\varphi':P\to Z$
so that $\rho_{i}\circ\varphi'=\pi_{i}$. We claim that $\varphi\circ\varphi'=1_{P}$
and $\varphi'\circ\varphi=1_{Z}$. Notice that $\pi_{i}\circ(\varphi\circ\varphi')=\rho_{i}\circ\varphi'=\pi_{i}$.
Again, by universal property of $P$ we find that $1_{P}$ is the
unique morphism which satisfies $\pi_{i}\circ1_{P}=\pi_{i}$for every
$1\leq i3$. This proves that $\varphi\circ\varphi'=1_{P}.$ Similarly,
by invoking universal property of $Z$, we find that $1_{Z}$ is the
unique morphism satisfiying $\rho_{i}\circ1_{Z}=\rho_{i}$ but this
means that it must be identical to the morphism $\varphi'\circ\varphi$
which also satisfies $\rho_{i}\circ(\varphi'\circ\varphi)=\pi_{i}\circ\varphi=\rho_{i}$.
This completes the proof that $P$ is unique up to isomorphism. $\hfill\blacksquare$
\begin{xca*}
\cite[Ch-1, Ex-5.10]{aluffi} Push the enveloope a little further
still, and define products and coproducts for families (i.e., indexed
sets) of objects of a category. 

Do these exist in $\mathsf{Set}.$ 
\end{xca*}
\begin{sol*}
Suppose we are given a family of objects $(A_{i})_{i\in I}$ in a
category $\catC$. An object $P=\prod_{I}A_{i}$ is said to be a product
of the $A_{i}$ together with a family of morphisms $(\pi_{i}:P\to A_{i})_{i\in I}$
if it satisfies the following universal property:

Given any object $Z\in\Obj(C)$ and a family of morphisms $(\rho_{i}:Z\to A_{i})_{i\in I}$
there exists a unique morphism $\varphi:Z\to P$ such that $\pi_{i}\circ\varphi=\rho_{i}$
for every $i\in I$. Such a $P$ is uniquely determined up to isomorphisms. 

I claim that given a family of sets $(A_{i})_{i\in I}$, their generalized
cartesian product 
\[
\prod_{i\in I}A_{i}\coloneqq\biggl\{ f:I\to\bigcup_{i\in I}A_{i}\mid f(i)\in A_{i}\text{ for every }i\in I\biggr\}
\]

is the product. Notice there are projections $\pi_{i}:\prod_{i\in I}A_{i}\to A_{i}$
for every $i\in I$ so that $\pi_{i}(f)=f(i)$. Given any object $C$
with a family of functions $(\rho_{i}:C\to A_{i})_{i\in I}$ there
exists a uniquely defined function $\varphi:C\to\prod_{i\in I}A_{i}$
such that for every $i\in I$, $\pi_{i}\circ\varphi=\rho_{i}$. Indeed,
simply define $\varphi(c)=f_{c}$ so that $f_{c}(i)=\rho_{i}(c)$.
That $\varphi$ is well defined as a function and satisfies the commutativity
relation for every $i\in I$ is obvious. We are left to show that
$\varphi$ is uniquely defined. For this, suppose $\varphi'$ is another
well defined map $C\to\prod_{i\in I}A_{i}$ so that $\pi_{i}\circ\varphi'=\rho_{i}$. 

We will show that $\varphi=\varphi'$. For this, take $c\in C,$ then
for all $i\in I$ we find that $(\varphi'(c))(i)=\pi_{i}(\varphi'(c))=\rho_{i}(c)=\pi_{i}(\varphi(c))=(\varphi(c))(i)$
and this completes the proof of uniqueness of $\varphi.$ $\hfill\blacksquare$
\end{sol*}
\begin{xca*}
\cite[Ch-1, Ex-5.11]{aluffi} Already did in a solution to some secion
1.1/1.2 problem. 
\end{xca*}
%
\begin{xca*}
\cite[Ch-1, Ex-5.12]{aluffi} Define the notions of \emph{fibered
products }and \emph{fibered coproducts,} as terminal objects of the
categories $\catC_{\alpha,\beta},\catC^{\alpha,\beta}$ considered
in Eg-3.10 by stating carefully the corresponding universal properties. 

As it happens, $\mathsf{Set}$ has both fibered products and fibered
coproducts. Define those objects 'concretely', in terms of naive set
theory. 
\end{xca*}
\begin{sol*}
TODO
\end{sol*}
\begin{thebibliography}{Aluffi, Algebra}
\bibitem[Aluffi, Algebra]{aluffi}

\end{thebibliography}

\end{document}
